\section*{第 \ref{ch:g6:classification-kinship} 章——分类与亲缘关系}

\begin{solution}{exo:g6:classification-kinship:1}
\omterm{def:g6:exploring-our-environment:components}{生物}所具有的、可供核查的一个\omterm{def:g6:classification-kinship:attribute}{特征}，陈述如下：
有\omterm{def:g4:classifying-animals:vertebrate}{脊柱}；有\omterm{def:g1:my-body:parts}{四肢}；有\omterm{def:g1:animals-around-us:covering}{皮毛}，并以\omterm{def:g2:animals-grow-up:born}{乳汁}哺育幼体
（或：有\omterm{def:g1:animals-around-us:covering}{羽毛}）。
\end{solution}

\begin{solution}{exo:g6:classification-kinship:2}
每个\omterm{def:g6:classification-kinship:attribute}{特征}的拥有者集合，都完整地落在另一个\omterm{def:g6:classification-kinship:attribute}{特征}的集合之内：
所有有\omterm{def:g1:animals-around-us:covering}{皮毛}的\omterm{def:g1:animals-around-us:animal}{动物}都有\omterm{def:g1:my-body:parts}{四肢}，所有有\omterm{def:g1:my-body:parts}{四肢}的\omterm{def:g1:animals-around-us:animal}{动物}都有\omterm{def:g4:classifying-animals:vertebrate}{脊柱}，
绝不会反过来，也绝不交叉——盒子套着盒子，
没有任何东西横跨两个盒子。
\end{solution}

\begin{solution}{exo:g6:classification-kinship:3}
因为它们继承自共同祖先：一个
\omterm{def:g6:classification-kinship:attribute}{特征}沿着家族谱系代代相传，因此拥有它的\omterm{def:g6:exploring-our-environment:components}{生物}都是该祖先的
后代——即一个家族。
\end{solution}

\begin{solution}{exo:g6:classification-kinship:4}
鳟鱼只越过了\omterm{def:g4:classifying-animals:vertebrate}{脊柱}这一个分叉点。
鲸越过了\omterm{def:g4:classifying-animals:vertebrate}{脊柱}、\omterm{def:g1:my-body:parts}{四肢}和\omterm{def:g1:animals-around-us:covering}{皮毛}加\omterm{def:g2:animals-grow-up:born}{乳汁}
这三个分叉点——全部都越过了。
\end{solution}

\begin{solution}{exo:g6:classification-kinship:5}
鲸与猫：它们共享\omterm{def:g1:animals-around-us:covering}{皮毛}加\omterm{def:g2:animals-grow-up:born}{乳汁}这一分叉，即最新画出的那一支；
鳟鱼则在最古老的分叉处就与它们分道扬镳。
\end{solution}

\begin{solution}{exo:g6:classification-kinship:6}
二者其实是同一信息的两种写法：一个盒子就是一个\omterm{def:g1:plants-around-us:tree}{分支}
连同它的全部小枝——分叉之后的一切都继承该分叉的\omterm{def:g6:classification-kinship:attribute}{特征}，
并落在该分叉所对应的盒子之中。
\end{solution}

\begin{solution}{exo:g6:classification-kinship:7}
\omterm{def:g1:animals-around-us:covering}{皮毛}加\omterm{def:g2:animals-grow-up:born}{乳汁}沿着一条家族谱系遗传，它标示着
\omterm{prop:g6:classification-kinship:kinship}{亲缘关系}。“像鱼一样的体形”则可能被毫无亲缘的\omterm{def:g6:exploring-our-environment:components}{生物}
因相同的水生生活方式而沾上，鲸就是例子：体形是借来的，
家族却在别处。分类划分的是家族，
因此它信任传家之宝，而不信任借来的外衣。
\end{solution}

\begin{solution}{exo:g6:classification-kinship:8}
鳟鱼：只有\omterm{def:g4:classifying-animals:vertebrate}{脊柱}。青蛙：\omterm{def:g4:classifying-animals:vertebrate}{脊柱}、\omterm{def:g1:my-body:parts}{四肢}。猫：
\omterm{def:g4:classifying-animals:vertebrate}{脊柱}、\omterm{def:g1:my-body:parts}{四肢}、\omterm{def:g1:animals-around-us:covering}{皮毛}加\omterm{def:g2:animals-grow-up:born}{乳汁}。乌鸫：\omterm{def:g4:classifying-animals:vertebrate}{脊柱}、
\omterm{def:g1:my-body:parts}{四肢}、\omterm{def:g1:animals-around-us:covering}{羽毛}。嵌套关系体现在勾号上：每一行的勾号都包含\omterm{def:g4:classifying-animals:vertebrate}{脊柱}；
凡标有\omterm{def:g1:animals-around-us:covering}{皮毛}加\omterm{def:g2:animals-grow-up:born}{乳汁}的行，都同时勾选了
\omterm{def:g1:my-body:parts}{四肢}。
\end{solution}

\begin{solution}{exo:g6:classification-kinship:9}
它越过了\omterm{def:g4:classifying-animals:vertebrate}{脊柱}和\omterm{def:g1:my-body:parts}{四肢}，并在自己的\omterm{def:g1:plants-around-us:tree}{分支}上离开，带着干燥、覆鳞的\omterm{def:g1:the-five-senses:sense}{皮肤}——
位于\omterm{def:g1:animals-around-us:covering}{皮毛}与\omterm{def:g1:animals-around-us:covering}{羽毛}分叉的旁边，而非越过
这两个分叉。
\end{solution}

\begin{solution}{exo:g6:classification-kinship:10}
哺乳标示\omterm{prop:g6:classification-kinship:kinship}{亲缘关系}：\omterm{def:g1:animals-around-us:covering}{皮毛}加\omterm{def:g2:animals-grow-up:born}{乳汁}是\omterm{prop:g3:sorting-living-things:groups}{哺乳动物}家族的传家之宝，
因此蝙蝠与猫同属一支。飞行是借来的生活方式——
翅膀在\omterm{prop:g3:sorting-living-things:groups}{鸟类}\omterm{def:g1:plants-around-us:tree}{分支}与蝙蝠\omterm{def:g1:plants-around-us:tree}{分支}上各自产生，
正如流线型体形在鲸与鳟鱼身上所见的那样，
同样是各自独立产生的。
\end{solution}

\begin{solution}{exo:g6:classification-kinship:11}
它属于\omterm{prop:g3:sorting-living-things:groups}{哺乳动物}。鸭嘴兽源自\omterm{prop:g3:sorting-living-things:groups}{哺乳动物}的共同祖先，并保有家族的
传家之宝（\omterm{def:g1:animals-around-us:covering}{皮毛}与\omterm{def:g2:animals-grow-up:born}{乳汁}）；其卵是一种古老的\omterm{def:g6:classification-kinship:attribute}{特征}，它的早期\omterm{def:g1:plants-around-us:tree}{分支}
保留了这一\omterm{def:g6:classification-kinship:attribute}{特征}，而其他\omterm{def:g1:plants-around-us:tree}{分支}则已失去。家族是祖先及其
\emph{全部}后代——包括古怪的成员在内；
这一组\omterm{def:g6:classification-kinship:attribute}{特征}，
（见\cref{exo:g4:classifying-animals:11}）已将其收入。
\end{solution}

\begin{solution}{exo:g6:classification-kinship:12}
若每一对\omterm{def:g1:plants-around-us:tree}{分支}都会相遇，则\omterm{def:g5:biodiversity:species}{物种}必然会在漫长的时间里
随世代更替而发生改变——新\omterm{def:g6:classification-kinship:attribute}{特征}不断产生并被遗传——
使得同一条祖先谱系能够分化出鳟鱼、猫与雏菊。
用一句话概括其含义：所有\omterm{def:g6:exploring-our-environment:components}{生物}都经过变化，从共同祖先
传承而来。需要索取的证据包括：变化的痕迹——
中间类型，古老岩层中保存的遗骸，
以及亲缘相似之处
恰好落在树形图所预言的位置上
（这属于以后某年的内容）。
\end{solution}

\begin{solution}{pb:g6:classification-kinship:1}
\textbf{1.} 鳟鱼：\omterm{def:g4:classifying-animals:vertebrate}{脊柱}。青蛙：\omterm{def:g4:classifying-animals:vertebrate}{脊柱}、\omterm{def:g1:my-body:parts}{四肢}。
蜥蜴：\omterm{def:g4:classifying-animals:vertebrate}{脊柱}、\omterm{def:g1:my-body:parts}{四肢}、覆鳞的干燥\omterm{def:g1:the-five-senses:sense}{皮肤}。鸽子：
\omterm{def:g4:classifying-animals:vertebrate}{脊柱}、\omterm{def:g1:my-body:parts}{四肢}、\omterm{def:g1:animals-around-us:covering}{羽毛}。老鼠：\omterm{def:g4:classifying-animals:vertebrate}{脊柱}、\omterm{def:g1:my-body:parts}{四肢}、
\omterm{def:g1:animals-around-us:covering}{皮毛}加\omterm{def:g2:animals-grow-up:born}{乳汁}。海豚：\omterm{def:g4:classifying-animals:vertebrate}{脊柱}、\omterm{def:g1:my-body:parts}{四肢}、
\omterm{def:g1:animals-around-us:covering}{皮毛}加\omterm{def:g2:animals-grow-up:born}{乳汁}（图注中关于\omterm{def:g2:animals-grow-up:born}{乳汁}的标注；以\omterm{def:g6:classification-kinship:attribute}{特征}为准，
而\omterm{def:g2:animals-grow-up:born}{乳汁}正是决定性的\omterm{def:g6:classification-kinship:attribute}{特征}）。

\textbf{2.} \omterm{def:g4:classifying-animals:vertebrate}{脊柱}——六者都是\omterm{def:g4:classifying-animals:vertebrate}{脊椎动物}。共享最广的
\omterm{def:g6:classification-kinship:attribute}{特征}标志着最古老的共同祖先：在该深度上，六者
同属一个家族。

\textbf{3.} 青蛙、蜥蜴、鸽子、老鼠、海豚。鸽子的翅膀与海豚的
鳍肢，本是同一对遗传下来的\omterm{def:g1:my-body:parts}{四肢}，只是各自\omterm{def:g4:adapted-to-their-world:adaptation}{适应}了不同的
生活方式，披上了不同的“工作装”；
系统树计入的是遗传关系，
而不是外表的装束。

\textbf{4.} 因为共享同一个居所，并不等于共享同一个
家族：“生活在水中”是一种\omterm{def:g2:where-animals-live:habitat}{栖息地}，可以被毫无亲缘的\omterm{def:g6:exploring-our-environment:components}{生物}借用，
并非身体祖先的\omterm{def:g6:classification-kinship:attribute}{特征}——这正是
即\cref{def:g3:sorting-living-things:criterion}中的旧规则，
如今附上了它的理由。

\textbf{5.} 树干 $\rightarrow$ \omterm{def:g4:classifying-animals:vertebrate}{脊柱}分叉（鳟鱼在此之后分出）
$\rightarrow$ \omterm{def:g1:my-body:parts}{四肢}分叉 $\rightarrow$ 三个
体被\omterm{def:g1:plants-around-us:tree}{分支}：干燥鳞片（蜥蜴）、\omterm{def:g1:animals-around-us:covering}{羽毛}（鸽子）、
\omterm{def:g1:animals-around-us:covering}{皮毛}与\omterm{def:g2:animals-grow-up:born}{乳汁}（小家鼠与海豚共处最新的分叉
上）；青蛙在\omterm{def:g1:my-body:parts}{四肢}分叉之后、各体被\omterm{def:g1:plants-around-us:tree}{分支}之前分出。

\textbf{6.} 小家鼠在“分类抽屉”中最近的亲属是
海豚（二者同属\omterm{def:g1:animals-around-us:covering}{皮毛}与\omterm{def:g2:animals-grow-up:born}{乳汁}分叉）。鳟鱼的最近亲属则是
其余所有\omterm{def:g6:exploring-our-environment:components}{生物}，彼此远近相同——它在第一个分叉处即已分出，
因此其余五种都是它同样疏远的亲属。

\textbf{7.} \omterm{prop:g3:sorting-living-things:groups}{两栖动物}——具有\omterm{def:g1:my-body:parts}{四肢}、\omterm{def:g1:the-five-senses:sense}{皮肤}裸露湿润，
没有属于自己的体被\omterm{def:g1:plants-around-us:tree}{分支}；其\omterm{def:g1:plants-around-us:tree}{分支}在\omterm{def:g1:my-body:parts}{四肢}分叉与
各体被\omterm{def:g1:plants-around-us:tree}{分支}之间分出。

\textbf{8.} 越过\omterm{def:g4:classifying-animals:vertebrate}{脊柱}与\omterm{def:g1:my-body:parts}{四肢}，进入干燥的
鳞片\omterm{def:g1:plants-around-us:tree}{分支}：\omterm{prop:g4:classifying-animals:classes}{爬行动物}——无需目睹即可归类，
因为决定归类的是\omterm{def:g6:classification-kinship:attribute}{特征}，而非标签。

\textbf{9.} 海豚与小家鼠关系更近：\omterm{def:g1:animals-around-us:covering}{皮毛}与\omterm{def:g2:animals-grow-up:born}{乳汁}是一个
遗传来的分叉。流线型体形则是水\omterm{def:g6:exploring-our-environment:components}{环境}的“统一制服”，
同样发给毫无亲缘的游泳者
（见\cref{exo:g4:adapted-to-their-world:11}）；若将其计入，
便会把借来的装束当成家族，打乱嵌套关系。

\textbf{10.} 一个共同祖先\omterm{def:g5:biodiversity:species}{物种}——一种早期的
\omterm{def:g1:my-body:parts}{四肢}\omterm{def:g4:classifying-animals:vertebrate}{脊椎动物}——青蛙、蜥蜴、鸽子、
小家鼠与海豚都由其演化而来，各自保留了
\omterm{def:g1:my-body:parts}{四肢}，只是穿着不同的装束。

\textbf{11.} 这六种\omterm{def:g6:exploring-our-environment:components}{生物}——以及画出它们的那一纲——
同属一个\omterm{def:g1:plants-around-us:tree}{分支}式的家族，
由共同祖先经演化而来，人类也处于
\omterm{def:g1:animals-around-us:covering}{皮毛}与\omterm{def:g2:animals-grow-up:born}{乳汁}\omterm{def:g1:plants-around-us:tree}{分支}上，与其余\omterm{def:g6:exploring-our-environment:components}{生物}并列。

\textbf{12.} 嵌套关系是\omterm{def:g6:exploring-our-environment:components}{生物}界的一个事实——\omterm{def:g6:classification-kinship:attribute}{特征}确实
呈盒中套盒的层级排列，这并非任何归档员
所规定。系统树解释了这一事实：来自共同
祖先的遗传\emph{必然}产生嵌套的\omterm{def:g6:classification-kinship:attribute}{特征}；
而单纯的归档，则会使这种恰好配合的整齐
成为一桩奇迹。
\end{solution}
